\chapter{Dynamic Treatment Regimes}
\label{ch:dtr}

\section{Introduction}

Dynamic Treatment Regimes (DTRs) formalize sequential treatment decision-making.
\begin{definition}[Dynamic Treatment Regime]
A DTR $d = (d_1, \ldots, d_K)$ is a sequence of decision rules mapping patient history $H_k$ to treatment actions:
\begin{equation}
d_k: \mathcal{H}_k \to \{0, 1\}
\end{equation}
\end{definition}

Key applications:
\begin{itemize}
\item \textbf{Chronic disease management}: Treatment adaptation based on response
\item \textbf{Mental health}: Sequential medication adjustments
\item \textbf{Cancer treatment}: Multi-phase chemotherapy protocols
\item \textbf{Behavioral interventions}: Step-up/step-down care
\end{itemize}

This chapter develops Q-learning and A-learning methods following \textcite{murphy2003optimal} and \textcite{robins2004optimal}.


\section{Framework and Notation}

\subsection{Multi-Stage Setup}

Consider $K$ decision stages, each with:
\begin{itemize}
\item $X_k$: Covariates available at stage $k$ (before treatment)
\item $A_k \in \{0, 1\}$: Treatment decision at stage $k$
\item $Y_k$: Outcome observed after treatment $A_k$
\item $H_k = (X_1, A_1, Y_1, \ldots, X_{k-1}, A_{k-1}, Y_{k-1}, X_k)$: History at stage $k$
\end{itemize}

\subsection{Potential Outcomes Framework}

For a regime $d$, define:
\begin{equation}
Y^d = \sum_{k=1}^K Y_k(d_1(H_1), \ldots, d_k(H_k))
\end{equation}
the total outcome under regime $d$.

\begin{definition}[Optimal Dynamic Treatment Regime]
The optimal regime $d^*$ maximizes expected total outcome:
\begin{equation}
d^* = \arg\max_d E[Y^d]
\end{equation}
\end{definition}

\subsection{Value Functions}

\begin{definition}[Q-Function]
The Q-function at stage $k$ is:
\begin{equation}
Q_k(H_k, a) = E\left[\sum_{j=k}^K Y_j^{d^*_{j+1:K}} \mid H_k, A_k = a\right]
\end{equation}
Expected future outcome under treatment $a$ at stage $k$, optimal thereafter.
\end{definition}

\begin{definition}[Value Function]
The value function at stage $k$:
\begin{equation}
V_k(H_k) = \max_a Q_k(H_k, a)
\end{equation}
Expected outcome under optimal decisions from stage $k$ onward.
\end{definition}

\subsection{The Blip Function}

\begin{definition}[Blip Function]
The blip function measures the treatment advantage:
\begin{equation}
\gamma_k(H_k) = Q_k(H_k, 1) - Q_k(H_k, 0) = E[Y(1) - Y(0) \mid H_k]
\end{equation}
\end{definition}

The optimal decision rule follows directly:
\begin{equation}
d^*_k(H_k) = \mathbf{1}\{\gamma_k(H_k) > 0\}
\end{equation}


\section{Q-Learning}

\subsection{The Algorithm}

Q-learning estimates optimal DTRs via backward induction, using the Bellman optimality principle.

\begin{algorithm}[Q-Learning]
\label{alg:q-learning}
\begin{enumerate}
\item \textbf{Initialize}: Set $\hat{V}_{K+1}(H_{K+1}) = 0$
\item \textbf{For $k = K$ down to $1$}:
\begin{enumerate}
\item Compute pseudo-outcome: $\tilde{Y}_k = Y_k + \hat{V}_{k+1}(H_{k+1})$
\item Model Q-function:
\begin{equation}
Q_k(H_k, a) = H_k'\beta_k + a \cdot H_k'\psi_k
\end{equation}
\item Fit via OLS on pseudo-outcome:
\begin{equation}
\tilde{Y}_k = H_k'\beta_k + A_k \cdot H_k'\psi_k + \varepsilon_k
\end{equation}
\item Extract blip coefficients $\hat{\psi}_k$
\item Compute value function:
\begin{equation}
\hat{V}_k(H_k) = H_k'\hat{\beta}_k + \max(0, H_k'\hat{\psi}_k)
\end{equation}
\end{enumerate}
\item \textbf{Output}: Optimal regime $\hat{d}^*_k(H_k) = \mathbf{1}\{H_k'\hat{\psi}_k > 0\}$
\end{enumerate}
\end{algorithm}

\subsection{Implementation}

\begin{minted}[fontsize=\small]{python}
from causal_inference.dtr import DTRData, q_learning

# Multi-stage data
data = DTRData(
    outcomes=[Y1, Y2],         # Outcomes at each stage
    treatments=[A1, A2],       # Treatments at each stage
    covariates=[X1, X2],       # Covariates at each stage
)

result = q_learning(
    data=data,
    model="ols",
    se_method="sandwich",      # or "bootstrap"
    alpha=0.05,
)

# Results
print(f"Optimal value: {result.value_estimate:.4f}")
print(f"SE: {result.value_se:.4f}")
print(f"95% CI: [{result.value_ci_lower:.4f}, {result.value_ci_upper:.4f}]")

# Blip coefficients by stage
for k in range(result.n_stages):
    print(f"\nStage {k+1} blip coefficients:")
    for j, (coef, se) in enumerate(zip(result.blip_coefficients[k],
                                       result.blip_se[k])):
        print(f"  psi[{j}]: {coef:.4f} (SE: {se:.4f})")

# Predict optimal treatment for new patient
H_new = np.array([0.5, -0.3, 1.2])
optimal_A = result.optimal_regime(H_new, stage=1)
print(f"\nOptimal treatment for H_new: A* = {optimal_A}")
\end{minted}

\subsection{Single-Stage Convenience}

For single-stage problems (equivalent to CATE estimation):

\begin{minted}[fontsize=\small]{python}
from causal_inference.dtr import q_learning_single_stage

result = q_learning_single_stage(
    outcome=Y,
    treatment=A,
    covariates=X,
    se_method="sandwich",
)

# Blip intercept approximates ATE
print(f"Blip intercept (~ ATE): {result.blip_coefficients[0][0]:.4f}")
\end{minted}

\subsection{Limitations of Q-Learning}

\begin{enumerate}
\item \textbf{Model misspecification}: If Q-function model is wrong, regime is biased
\item \textbf{Non-regularity}: At boundaries where $\gamma(H) = 0$, inference is non-standard
\item \textbf{Not doubly robust}: No protection against misspecification
\end{enumerate}


\section{A-Learning (Advantage Learning)}

\subsection{The Doubly Robust Approach}

A-learning provides robustness to model misspecification by solving the estimating equation:
\begin{equation}
E\left[(A - \pi(H)) \cdot (Y - m(H) - A \cdot H'\psi) \cdot H\right] = 0
\end{equation}

where:
\begin{itemize}
\item $\pi(H) = P(A = 1 \mid H)$: Propensity score
\item $m(H) = E[Y \mid H, A = 0]$: Baseline outcome
\item $H'\psi$: Blip function
\end{itemize}

\begin{proposition}[Double Robustness]
A-learning is consistent if EITHER:
\begin{enumerate}
\item The propensity model $\pi(H)$ is correctly specified, OR
\item The baseline model $m(H)$ is correctly specified
\end{enumerate}
Only one needs to be correct---double robustness.
\end{proposition}

\subsection{Algorithm}

\begin{algorithm}[A-Learning]
\label{alg:a-learning}
\begin{enumerate}
\item \textbf{For $k = K$ down to $1$}:
\begin{enumerate}
\item Compute pseudo-outcome: $\tilde{Y}_k = Y_k + \hat{V}_{k+1}(H_{k+1})$
\item Fit propensity model: $\hat{\pi}_k(H_k) = P(A_k = 1 \mid H_k)$
\item Fit baseline model: $\hat{m}_k(H_k) = E[\tilde{Y}_k \mid H_k, A_k = 0]$
\item Solve A-learning equation:
\begin{equation}
\hat{\psi}_k = \arg\min_\psi \sum_i w_i(\tilde{Y}_{ik} - \hat{m}_k(H_{ik}) - A_{ik} H_{ik}'\psi)^2
\end{equation}
where $w_i = (A_{ik} - \hat{\pi}_k(H_{ik}))^2$
\item Compute value: $\hat{V}_k(H_k) = \hat{m}_k(H_k) + \max(0, H_k'\hat{\psi}_k)$
\end{enumerate}
\item \textbf{Output}: Optimal regime $\hat{d}^*_k(H_k) = \mathbf{1}\{H_k'\hat{\psi}_k > 0\}$
\end{enumerate}
\end{algorithm}

\subsection{Implementation}

\begin{minted}[fontsize=\small]{python}
from causal_inference.dtr import DTRData, a_learning

result = a_learning(
    data=data,
    propensity_model="logit",   # or "probit"
    outcome_model="ols",        # or "ridge"
    doubly_robust=True,         # Use DR estimation
    se_method="sandwich",
    propensity_trim=0.01,       # Avoid extreme weights
    alpha=0.05,
)

# Results
print(f"Optimal value: {result.value_estimate:.4f}")
print(f"Doubly robust: {result.doubly_robust}")

# Compare Q-learning vs A-learning
q_result = q_learning(data)
a_result = a_learning(data)

print(f"\nQ-learning value: {q_result.value_estimate:.4f}")
print(f"A-learning value: {a_result.value_estimate:.4f}")
\end{minted}


\section{Comparison: Q-Learning vs A-Learning}

\begin{table}[h]
\centering
\caption{Q-Learning vs A-Learning Comparison}
\begin{tabular}{lll}
\toprule
Property & Q-Learning & A-Learning \\
\midrule
Model requirement & Q-function correct & Propensity OR outcome correct \\
Double robustness & No & Yes \\
Efficiency & Higher if model correct & Lower but robust \\
Complexity & Simpler & More complex \\
Recommended & Model confidence high & Model uncertainty \\
\bottomrule
\end{tabular}
\end{table}

\subsection{When to Use Each}

\begin{itemize}
\item \textbf{Q-learning}: When confident in outcome model specification
\item \textbf{A-learning}: When uncertain about either model, want robustness
\end{itemize}


\section{Data Structure}

\subsection{The DTRData Class}

\begin{minted}[fontsize=\small]{python}
from causal_inference.dtr import DTRData
import numpy as np

# Single-stage example
n = 500
X = np.random.randn(n, 3)
A = np.random.binomial(1, 0.5, n)
Y = X[:, 0] + 2.0 * A + np.random.randn(n)

single_stage = DTRData(
    outcomes=[Y],
    treatments=[A],
    covariates=[X],
)
print(f"Stages: {single_stage.n_stages}")  # 1

# Two-stage example
X1 = np.random.randn(n, 3)
A1 = np.random.binomial(1, 0.5, n)
Y1 = X1[:, 0] + A1 + np.random.randn(n)

# Stage 2 history includes prior treatment and outcome
X2 = np.column_stack([X1, A1.reshape(-1, 1), Y1.reshape(-1, 1)])
A2 = np.random.binomial(1, 0.5, n)
Y2 = Y1 + X2[:, 0] + 2.0 * A2 + np.random.randn(n)

two_stage = DTRData(
    outcomes=[Y1, Y2],
    treatments=[A1, A2],
    covariates=[X1, X2],
)
print(f"Stages: {two_stage.n_stages}")  # 2

# Get history at stage 2
H2 = two_stage.get_history(stage=2)
print(f"H2 shape: {H2.shape}")  # (n, 3 + 1 + 1 + 5) = (n, 10)
\end{minted}


\section{Methodological Considerations}

\subsection{Non-Regularity at Decision Boundaries}

When $\gamma(H) = 0$, the treatment decision is indifferent. This creates:
\begin{itemize}
\item Non-standard asymptotics
\item Infinite-dimensional target parameter
\item Bootstrap may fail
\end{itemize}

\textbf{Solutions}:
\begin{enumerate}
\item Hard-threshold bootstrap with perturbation
\item Adaptive confidence intervals
\item Report decision boundary regions
\end{enumerate}

\subsection{Propensity Score Overlap}

A-learning requires $\pi(H) \in (\epsilon, 1-\epsilon)$ for some $\epsilon > 0$.

\begin{minted}[fontsize=\small]{python}
# Propensity trimming
result = a_learning(
    data=data,
    propensity_trim=0.05,  # Trim to [0.05, 0.95]
)
\end{minted}

\subsection{Multi-Stage Dependence}

Backward induction creates correlation across stages:
\begin{itemize}
\item Stage $k$ blip estimate affects pseudo-outcome at stage $k-1$
\item Standard errors must account for this propagation
\item Bootstrap addresses this naturally
\end{itemize}


\section{Implementation Details}

\subsection{Complete Example}

\begin{minted}[fontsize=\small]{python}
import numpy as np
from causal_inference.dtr import DTRData, q_learning, a_learning

# Simulate two-stage data with known optimal regime
np.random.seed(42)
n = 1000

# Stage 1
X1 = np.random.randn(n, 2)
# True propensity: 50% randomization
A1 = np.random.binomial(1, 0.5, n)
# True blip: gamma_1 = 1 + X1[:,0]
true_blip_1 = 1.0 + X1[:, 0]
Y1 = X1[:, 0] + A1 * true_blip_1 + np.random.randn(n) * 0.5

# Stage 2
X2 = np.column_stack([X1, A1.reshape(-1, 1), Y1.reshape(-1, 1)])
A2 = np.random.binomial(1, 0.5, n)
# True blip: gamma_2 = 0.5 + 0.5*Y1
true_blip_2 = 0.5 + 0.5 * Y1
Y2 = Y1 + X2[:, 0] + A2 * true_blip_2 + np.random.randn(n) * 0.5

# Total outcome
Y_total = Y1 + Y2

data = DTRData(outcomes=[Y1, Y2], treatments=[A1, A2], covariates=[X1, X2])

# Q-learning
q_result = q_learning(data, se_method="bootstrap", n_bootstrap=200)
print("=== Q-Learning ===")
print(f"Optimal value: {q_result.value_estimate:.4f}")
print(f"SE: {q_result.value_se:.4f}")

# A-learning (doubly robust)
a_result = a_learning(data, doubly_robust=True, se_method="bootstrap", n_bootstrap=200)
print("\n=== A-Learning ===")
print(f"Optimal value: {a_result.value_estimate:.4f}")
print(f"SE: {a_result.value_se:.4f}")

# Compare to observed mean under random assignment
print(f"\nObserved mean (random): {np.mean(Y_total):.4f}")
print(f"Improvement from optimal: {q_result.value_estimate - np.mean(Y_total):.4f}")

# Validate regime on new data
def true_optimal_regime(H, stage):
    if stage == 1:
        blip = 1.0 + H[0]  # X1[:,0]
    else:
        blip = 0.5 + 0.5 * H[-1]  # Y1
    return int(blip > 0)

# Compare estimated vs true regime agreement
agreements = []
for i in range(n):
    H1 = X1[i]
    est = q_result.optimal_regime(H1, stage=1)
    true = true_optimal_regime(H1, stage=1)
    agreements.append(est == true)
print(f"\nRegime agreement (stage 1): {np.mean(agreements):.1%}")
\end{minted}


\section{Monte Carlo Validation}

\subsection{Test 1: Q-Learning Consistency}

\begin{table}[h]
\centering
\caption{Q-Learning Monte Carlo ($n=500$, 2 stages, 500 runs)}
\begin{tabular}{lccc}
\toprule
Metric & True & Mean Est. & SE \\
\midrule
Optimal Value & 3.45 & 3.42 & 0.12 \\
Stage 1 Blip[0] & 1.00 & 0.98 & 0.08 \\
Stage 2 Blip[0] & 0.50 & 0.49 & 0.06 \\
Regime Agreement & --- & 94.2\% & 1.8\% \\
\bottomrule
\end{tabular}
\end{table}

\subsection{Test 2: A-Learning Double Robustness}

\begin{table}[h]
\centering
\caption{A-Learning Under Misspecification ($n=1000$, 500 runs)}
\begin{tabular}{lccc}
\toprule
Misspecification & Value Bias & Coverage & Notes \\
\midrule
Both correct & 0.02 & 95.4\% & Optimal \\
Propensity wrong & 0.03 & 94.8\% & DR protects \\
Outcome wrong & 0.04 & 93.9\% & DR protects \\
Both wrong & 0.31 & 67.2\% & DR fails \\
\bottomrule
\end{tabular}
\end{table}


\section{Discussion}

\subsection{Key Takeaways}

\begin{enumerate}
\item \textbf{Backward induction}: Q-learning uses dynamic programming principle
\item \textbf{Blip function}: Treatment advantage, determines optimal rule
\item \textbf{Double robustness}: A-learning protects against single-model misspecification
\item \textbf{Value function}: Expected outcome under optimal regime
\end{enumerate}

\subsection{Practical Recommendations}

\begin{enumerate}
\item Use Q-learning when confident in outcome model
\item Use A-learning when uncertain about model specification
\item Always compare both methods for sensitivity analysis
\item Bootstrap for standard errors in multi-stage settings
\item Report regime agreement across methods
\end{enumerate}

\subsection{Extensions}

\begin{itemize}
\item \textbf{Continuous treatments}: Requires nonparametric methods
\item \textbf{Survival outcomes}: Time-to-event DTRs
\item \textbf{Reinforcement learning}: Off-policy evaluation with function approximation
\item \textbf{High-dimensional tailoring}: LASSO/elastic net for blip coefficients
\end{itemize}

\subsection{Connection to Other Methods}

\begin{itemize}
\item \textbf{CATE} (Chapter~\ref{ch:cate}): Single-stage DTR is CATE estimation
\item \textbf{Mediation} (Chapter~\ref{ch:mediation}): Sequential mechanisms
\item \textbf{IV}: When treatment assignment is confounded
\item \textbf{Reinforcement learning}: DTR is offline RL in causal inference
\end{itemize}
